% ============================================================
%  1. Introduction  (~0.5 page)
% ============================================================
\section{Introduction}

\subsection{Contexte}
Ce projet a pour but de concevoir, à l'aide du langage \textbf{Java} 
et du générateur d'analyseurs \textbf{JavaCC}, un Système de Gestion 
de Base de Données (SGBD) relationnel. Il s'agit d'un interpréteur 
de commandes SQL permettant à un utilisateur de créer et manipuler 
une base de données directement depuis un terminal.

Le système repose sur une architecture orientée objet structurée en 
plusieurs packages : modèle de données, expressions logiques, 
commandes, analyse syntaxique et exécution.

\subsection{Fonctionnalités implémentées}
Le SGBD supporte les commandes SQL suivantes :

\begin{itemize}
    \item \texttt{CREATE TABLE} - création d'une table avec ses attributs et leurs types
    \item \texttt{INSERT INTO} -insertion d'un tuple dans une table
    \item \texttt{SELECT} - consultation des données avec projection, 
          filtrage (\texttt{WHERE}) et jointures multi-tables
    \item \texttt{DELETE FROM} - suppression sélective de tuples selon une condition
    \item \texttt{DROP TABLE} -suppression définitive d'une table
    \item \texttt{TABLES} - affichage de la liste des tables existantes
    \item \texttt{EXIT} / \texttt{QUIT} - fermeture de l'application 
          avec sauvegarde automatique
\end{itemize}

Les types de données supportés sont \texttt{INT}, \texttt{VARCHAR} et \texttt{SERIAL} (entier auto-incrémenté).